\begin{abstract}
Organisations run several business processes side by side and are measured on
organisation-level indicators. Reinforcement learning (RL) for prescriptive process
management usually treats such an organisation as one monolithic decision problem:
every decision is credited with the reward noise of every process, the only way to
wait for a better resource is to freeze the whole organisation, and every decision
reasons over the entire organisation. We observe that the process model already
certifies which parts of the organisation are independent: the connected components
of its Action-Evolution Petri Net. We use this independence certificate at three points
of the RL pipeline: to credit each decision only with the rewards of its own process
(\nfgae), to let a single process postpone its decisions without freezing the others
(component-scoped postponement), and to let each decision observe only its own process
(component turns). Each use is exact where the model certifies independence and
reduces to standard PPO on a single process. On an insurer benchmark with claims,
underwriting and complaints processes, standard PPO degrades as the number of
concurrent processes grows, whereas the independence-aware pipeline does not; it
learns to wait exactly where waiting pays, and only together with factored credit,
outperforming a waiting heuristic, at a compute cost linear in the number of processes.
With six concurrent processes, PPO ends at 0.75 of the heuristic's return and factored
credit at 1.07, ahead on 20 of 20 seeds; where waiting pays, factored credit with
component-scoped postponement is the only arm that beats the waiting heuristic.

\keywords{Prescriptive process management \and Reinforcement learning \and
Petri nets \and Resource allocation \and Credit assignment}
\end{abstract}
